% -----------------------------
% 執筆にあたっての留意点
% -----------------------------
% 結果に意味を与える章であり，イントロダクションに次いで重要な章です．考察の章の役割は，結果の章で示した実験結果，調査結果を踏まえて仮説の検証を行い，結論にいたるまでの議論を展開することにあります．結果の章では実験や調査で得られた結果を淡々と報告すると述べましたが，考察の章に入ってようやく結果の解釈を行い，著者の「意見」を述べることができます．
%
% 考察の章では以下の項目について記します．
%
% * 研究目的の振り返り
% * 結果の解釈と仮説の立証結果
% * 提案手法，研究手法の制約条件，限界
% * 提案手法がうまく機能した理由，うまく機能しなかった理由の考察
% * 結果から導かれる提言
%
% 考察は，論文の中で最も執筆に神経を使う章です．特に文章の論理展開には気をつけてください．実験・調査結果を踏まえて何らかの主張を行うときは，論理的な飛躍がないように注意しましょう．何らかの提言を行う場合は，立証された仮説や分析結果に関連している必要があります．実験・調査結果をよりどころとしない主張はしてはいけません．
%
%
% -----------------------------
% 補足（参考：how-to-xx/how-to-write-a-paper/1-学術論文の書き方.md 3.5節，
%              how-to-xx/how-to-write-a-paper/3-国際会議論文の傾向と対策.md 2章）
% -----------------------------
%
% ■ 結果の章の要約を繰り返すだけの章にしてはいけません
% 「図Xでは〜だった」と再掲するのではなく，「なぜそうなったのか」「それは何を意味するのか」を書きます．
% 結果の再掲は，議論を成立させるために必要な最小限にとどめてください．
%
% ■ 主張の一つ一つに，対応する結果があるかを確認します
% 考察で述べる主張は，結果の章のどの数値・どの図表に支えられているかを言えなければなりません．
% 「〜だと思われる」「〜であろう」と書きたくなったら，その根拠が結果の章にあるかを確認してください．
% なければ，それは考察ではなく単なる感想です．
%
% ■ 限界（Limitations）は，隠すのではなく自分から書きます
% 実験協力者の偏り，タスクの人工性，一般化できる範囲などを，自分から具体的に述べてください．
% 査読者が気づく限界を著者が書いていない場合，「気づいていない」か「隠している」と受け取られます．
% 自分の研究の限界を正確に把握できていることは，研究者としての能力の証明になります．
% なお，「限界を書く」ことと「自信のなさを表明する」ことは違います．
% 「〜という限界はあるが，本研究の主張は〜の範囲では成立する」と，主張の有効範囲を示す書き方をしてください．
%
% ■ HCI系の会議に投稿する場合
% 得られた知見を「今後の設計に対する示唆（design implications）」の形でまとめるのが定番の型です．
% 「本研究の結果は，〜のようなインタフェースを設計する際に〜すべきことを示唆している」といった形で書きます．

% ----------------------
% 考察
% ----------------------
\section{考察}
